% ---------------------------------------------------------------------
% 2.2 Fundamentos neurofisiológicos
% ---------------------------------------------------------------------
\section[FUNDAMENTOS NEUROFISIOLÓGICOS]{Fundamentos neurofisiológicos}
\label{sec:neuro}

Para entender de dónde proviene la señal que se busca clasificar es necesario revisar algunos
conceptos de la fisiología del sistema nervioso, en particular la forma en que las neuronas generan
y transmiten actividad eléctrica y la organización de la corteza cerebral, que es la región desde
la cual se origina la mayor parte de la actividad registrada en la superficie de la cabeza.

\subsection{La neurona y el impulso nervioso}

La neurona es la unidad funcional del sistema nervioso y se compone de un cuerpo celular o soma,
de un conjunto de prolongaciones cortas llamadas dendritas, que reciben las señales de otras
células, y de una prolongación larga llamada axón, que conduce la señal hacia las neuronas
siguientes. La membrana neuronal mantiene en reposo una diferencia de potencial entre el interior y
el exterior de la célula, cercana a $-70$~mV, gracias a la distribución desigual de iones de sodio y
potasio; cuando un estímulo despolariza la membrana por encima de un umbral, se abren canales
iónicos dependientes de voltaje y se genera un potencial de acción que se propaga a lo largo del
axón \cite{guyton2016}. Al llegar a la terminal axónica, la señal se transmite a la neurona
siguiente mediante la sinapsis química, en la que la liberación de neurotransmisores produce en la
célula receptora potenciales postsinápticos excitadores o inhibidores, cuya duración, de decenas
de milisegundos, es considerablemente mayor que la de un potencial de acción, que dura
aproximadamente un milisegundo \cite{ganong2010}.

\subsection{La corteza cerebral y sus lóbulos}

El cerebro es la porción más voluminosa del encéfalo y su capa más externa, la corteza cerebral, es
la responsable de funciones como la percepción, el lenguaje, la memoria y el control voluntario del
movimiento. Cada hemisferio se divide en cuatro lóbulos principales, cuyo nombre corresponde al
hueso del cráneo que los cubre: el lóbulo frontal, relacionado con la planeación, la toma de
decisiones y el control motor; el lóbulo parietal, que integra la información somatosensorial y
participa en los procesos de atención espacial; el lóbulo temporal, asociado a la audición y a la
memoria; y el lóbulo occipital, que contiene la corteza visual primaria \cite{snell2014}. Esta
división resulta relevante para el presente trabajo porque la nomenclatura de los electrodos del
sistema internacional 10-20, descrita más adelante, toma precisamente la inicial de cada lóbulo, y
porque la componente P300 se registra con mayor amplitud sobre las regiones centro-parietales de la
línea media \cite{polich2007}.

\subsection{Origen de la actividad eléctrica medible}

La descarga de una sola neurona no puede detectarse desde la superficie de la cabeza, pues para
que un sensor colocado sobre el cuero cabelludo registre actividad es necesario que miles o incluso
millones de neuronas se activen de forma sincrónica \cite{guyton2016}. La mayor parte de esta
señal proviene de la suma de los potenciales postsinápticos de las neuronas piramidales de la
corteza, que se encuentran alineadas de manera perpendicular a la superficie cortical y cuya
actividad, al ocurrir en sincronía, genera dipolos de corriente que se suman en lugar de
cancelarse \cite{niedermeyer2005}. Por esta razón, la intensidad de las ondas registradas depende
más del número de neuronas que disparan de manera sincronizada que del nivel total de actividad
eléctrica del encéfalo, y señales potentes pero asincrónicas pueden anularse entre sí debido a su
polaridad opuesta \cite{guyton2016}.

\subsection{Atención y procesamiento de estímulos visuales}

El paradigma que se utiliza en este trabajo aprovecha la manera en que el cerebro responde a los
estímulos que resultan relevantes para una tarea. Cuando una persona observa una secuencia de
estímulos visuales y tiene que identificar uno en particular, la información sensorial pasa por una
etapa inicial de procesamiento en la corteza visual y posteriormente por un proceso de comparación,
dirigido por la atención, contra la representación del contexto que se mantiene en la memoria de
trabajo \cite{polich2007}. Si el estímulo coincide con aquel al que la persona presta atención, se
produce una actualización de dicha representación que se refleja en el EEG como una deflexión
positiva alrededor de 300~ms después del estímulo, fenómeno que da origen a la componente P300
descrita en la Sección~\ref{sec:erp}.
